\documentclass{article}
\usepackage{lipsum}

\usepackage[LGR, T1]{fontenc} %για γλώσσα
\usepackage[greek,english]{babel}
\usepackage[utf8]{inputenc}

\usepackage{graphicx} %για εικόνες
\usepackage{float} 
\usepackage{wrapfig} 
\usepackage{color}   %May be necessary if you want to color links
\usepackage{textcomp}
\usepackage{xcolor}

\usepackage{caption}
\usepackage{hyperref}

\begin{document}
\selectlanguage{Greek}
\begin{titlepage}

    \begin{figure}
        \centering
        \includegraphics[scale=0.4]{Upatras}
    \end{figure}

    \begin{center}
        \line(1,0){300}\\
        [0.25in]
        \huge{\bfseries \selectlanguage{English}Project 2021-2022}\\

        \line(1,0){200}\\
        \begin{figure}[H]
            \centering
            \includegraphics[scale=2.2]{logo}\\
        \end{figure}
    \end{center}

    \begin{flushright}
        \textsc{\\}
        \textsc{\\}

        \textsc{\large \\ Τριανταφυλλόπουλος Παναγιώτης\\ΑΜ 1054367}
        \textsc{\large \\ Τμήμα Μηχ. Η/Υ \& Πληροφορικής}
        \textsc{\large \\ Αναπαράσταση Γνώσης στον Παγκόσμιο Ιστό}
        \textsc{\large \\ \selectlanguage{English}Project \selectlanguage{Greek} 2021-2022}
    \end{flushright}

\end{titlepage}

\tableofcontents
\thispagestyle{empty}
\cleardoublepage

\setcounter{page}{1}

\selectlanguage{greek}
\section{Εισαγωγή}
\addcontentsline{toc}{section}{Εισαγωγή}

Καλείστε να συμμετάσχετε στη συλλογή δεδομένων στα πλαίσια της διπλωματικής εργασίας του προπτυχιακού φοιτητή Τριανταφυλλόπουλου Παναγιώτη, η οποία έχει σκοπό τη δημιουργία ενός συστήματος συστάσεων ιστοσελίδων. Για τη συλλογή δεδομένων θα χρειαστείτε έναν φυλλομετρητή (selectlanguage{english}browser\selectlanguage{greek} - κατά προτίμηση selectlanguage{english}Google Chrome\selectlanguage{greek} ή selectlanguage{english}Firefox\selectlanguage{greek}) και να εγκαταστήσετε και να χρησιμοποιήσετε μία επέκταση φυλλομετρητή (\selectlanguage{english}browser extension\selectlanguage{greek}) που έχει δημιουργηθεί από τον φοιτητή.


\selectlanguage{english}
\section{Disclaimer}
\selectlanguage{greek} Αυτό το ερευνητικό έργο θα χρησιμοποιήσει έναν τυχαία δημιουργημένο κωδικό αναγνώρισης για να διασφαλίσει την ανωνυμία όλων των συμμετεχόντων. Δε θα συλλεχθούν ούτε θα αποθηκευθούν προσωπικά στοιχεία ταυτοποίησης και δεν θα γίνει προσπάθεια σύνδεσης οποιουδήποτε συγκεκριμένου συμμετέχοντος με τις απαντήσεις ή τα δεδομένα του. Όλα τα δεδομένα που συλλέγονται κατά τη διάρκεια αυτής της μελέτης θα διατηρηθούν αυστηρά απόρρητα και θα χρησιμοποιηθούν αποκλειστικά για ερευνητικούς σκοπούς. Η συμμετοχή σε αυτή τη μελέτη είναι εντελώς εθελοντική και οι συμμετέχοντες μπορούν να αποσυρθούν ανά πάσα στιγμή χωρίς κυρώσεις ή απώλεια ωφελημάτων που διαφορετικά δικαιούνται.

\section{Οδηγίες Εγκατάστασης}

\section{Οδιγίες Χρήσης}


\end{document}
